\chapter*{全书导读}
\addcontentsline{toc}{chapter}{全书导读}
\markboth{全书导读}{全书导读}

\section*{一、这本书是什么}

2022 年 7 月，IHES 举办了为期三周的 Langlands 纲领暑期学校，纪念 1977 年 Corvallis 会议 45 周年。会议录分三卷出版，本卷是第二卷，主题为\textbf{Langlands 对应的几何化}（geometrization）。六篇文章覆盖了当下几何 Langlands 最活跃的几条线索：局部 Langlands 参数的模空间、其上的凝聚层、曲线的模叠、$p$-进 Langlands 的范畴化、凝聚型与可构造型局部 Langlands 对应的关系，以及导出方法。

如果把经典 Langlands 纲领理解为「$L$-函数与自守表示的对应」，那么几何化讨论的是更细的问题：\emph{当参数本身构成一个几何对象（叠）时，表示一侧应当是什么？}答案的候选者包括凝聚层、可构造层、导出范畴、以及各种「族」中的对应。六篇文章从不同角度逼近这一问题。

\section*{二、六篇文章速览}

\paragraph{第 1 篇（Dat）局部 Langlands 参数的模空间。}
Dat 的文章是全书最好的入口。他先解释为什么要把 Langlands 参数（$W_F \to \hat G$ 的 1-余循环，模去 $\hat G$-共轭）做成一个代数叠 $Z^1/\hat G$，然后处理一个技术核心：如何让「连续性」在 $\ell$-进与 $p$-进系数下都有意义（用 $\hat G$ 的 $\mathbb Z_\ell$-模型与 Weil 群的拓扑生成元）。第 3 节用「驯顺参数」这个玩具模型把整个构造走一遍，第 4 节把一般情形归约到驯顺情形并计算连通分支，第 5 节讨论粗模空间及其在「族中的 LLC」中的应用。读完全书其他文章后再回来读第 5 节，会有第二次收获。

\paragraph{第 2 篇（Zhu）参数叠上的凝聚层。}
Zhu 的文章最长也最系统（85 页），实际上是一部小型专著。它把 $\mathrm{Qcoh}(\mathcal L_G)$ 当作局部 Langlands 对应的「Betti 侧」替身，构造了从光滑表示范畴到参数叠上凝聚层的函子，并证明它在许多情形下是范畴等价。文中的关键工具是「参数叠上的导出代数几何」——参数叠不是拟紧的，需要用 $\infty$-范畴与导出完备化的语言。第 6 节之后的算例（环面、$\mathrm{GL}_2$、分裂群）值得逐一验算。

\paragraph{第 3 篇（Fargues–Scholze）曲线上的向量丛模叠。}
这是 Fargues–Scholze 系列工作的综述，把局部 Langlands 对应重新表述为曲线 $X$ 上的 $\mathrm{Bun}_G$ 与参数叠之间的几何对应。核心猜想断言 \textbf{局部 Langlands 对应的几何化等价于 }$\mathrm{Bun}_G$ \textbf{上的凝聚层与参数叠上的凝聚层之间的等价}，并用曲线的几何（Harder–Narasimhan 约化、仿射 Grassmann 的施坦结构）给出了一条证明路线。注意本文与第 2 篇的语言差异：Zhu 用参数叠，Fargues–Scholze 用「钻石」（diamond）与 $v$-叠。

\paragraph{第 4 篇（Emerton–Gee–Hellmann）$p$-进 Langlands 范畴化导论。}
全书最长的一篇（253 页），也是最像教科书的一篇。它从 $\mathrm{GL}_2(\mathbb Q_p)$ 的 Banach 表示出发，逐步建立「$p$-进 Langlands 对应的范畴化」：以 $\mathrm{Gal}(\overline{\mathbb Q}_p/\mathbb Q_p)$ 的表示范畴为模型，构造 $\mathrm{GL}_n(\mathbb Q_p)$ 的「解析表示」范畴与其间的等价。六个附录（范畴论速成、IndCoh 与 ProCoh、局部化与 Morita 理论、拓扑阿贝尔群、$p$-进解析群的表示、群表示的导出模叠）自成体系，可以单独当参考书读。若只读一篇，读这一篇；若要动手算，先读第 1 篇。

\paragraph{第 5 篇（Ben-Zvi–Chen–Helm–Nadler）凝聚型与可构造型之间。}
本文篇幅最短，问题最尖锐：局部 Langlands 对应的两种精细形式——以 Vogan、Soergel 为代表的\emph{可构造层}版本（Kazhdan–Lusztig 风格，固定中心特征），与以 Fargues、Zhu、Hellmann 为代表的\emph{凝聚层}版本（几何 Langlands 风格，中心特征可变、参数空间更大且允许 Frobenius 的单幂部分）——究竟是什么关系？作者提出一个统一的机制：用导出代数几何中的\textbf{圆周作用}（circle action）与循环同调来内插两种设定。非阿基米德情形下，他们说明圆周作用如何联系仿射 Hecke 代数与 $\mathrm{GL}_n$ 全体光滑表示的凝聚实现与可构造实现；阿基米德情形下，他们从 twistor 直线上的驯顺几何 Langlands 猜想（Fargues 猜想的实形式）出发，导出 Adams–Barbasch–Vogan 与 Soergel 形式的可构造对应。技术核心是 loop space 的 Koszul 对偶。

\paragraph{第 6 篇（Feng–Harris）Langlands 对应中的导出结构。}
Feng–Harris 的综述聚焦「导出」这一层：导出 Galois 形变环（Galatius–Venkatesh）、导出 Hecke 代数（Ollivier–Schneider、Venkatesh 两派工作）、导出 Hitchin 叠（Feng–Yun–Zhang）、导出特殊圈与高阶 theta 函数。他们的论点是：Galois 侧与自守侧都存在自然的导出结构，而对应在导出层面才真正自然；文中还讨论局部对称空间的上同调、Venkatesh 的动机猜想，并列出一批开放问题。附录 A 是单纯交换环速成，读懂正文前建议先读附录。

\section*{三、阅读路线建议}

\begin{itemize}[leftmargin=2em,itemsep=0.25em]
\item \textbf{初学（有 $p$-进表示基础）}：第 1 篇 → 第 4 篇正文（跳过附录，遇范畴论语言再回查）→ 第 3 篇。
\item \textbf{想快速了解「几何化是什么」}：第 1 篇前两节 + 第 5 篇引言 + 本导读。约一天。
\item \textbf{做凝聚侧}：第 1 篇 → 第 2 篇 → 第 5 篇 → 第 6 篇。
\item \textbf{做 $p$-进侧}：第 4 篇（尤其附录 B、C、E、F）→ 第 3 篇 → 第 6 篇。
\end{itemize}

\section*{四、预备知识}

读懂全书需要：局部域与 Galois 上同调、约化群的表示（Bernstein 中心、Whittaker 模型）、$\ell$-进层与六函子、导出范畴与 $\infty$-范畴的基本语言。不需要先读完全部预备文献——第 1 篇与第 4 篇的引言都刻意写得自足，读到不懂的术语再查更有效率。

\section*{五、与另外两卷的关系}

第一卷讲表示、内窥、志村簇与 shtuka 的上同调；第三卷讲相对 Langlands（周期、theta 对应、算术版本）。本卷处于中间：它假定读者知道「对应应当存在」，转而追问「对应在什么几何舞台上发生」。三卷合读，才能看到 Langlands 纲领今日的全貌。

\section*{六、如何使用本译本}

译文保留了底本的全部编号与页锚（切口侧灰色数字）。公式均逐页核对过底本图像，底本笔误以脚注标出。术语以 \texttt{glossary.tsv} 为准；遇到译名不一致，以术语表为最终依据。译者注只补充背景与澄清，不改变论证。若要引用，请引用原书页码（页锚）而非本译本页码。
